\chapter{量子谐振子}\label{ch:b3:harmonic-oscillator}

自然界中几乎没有任何东西严格是谐振子，而几乎一切都
近似为谐振子：任何被扰动而离开稳定平衡位置的系统——
分子的化学键、晶体中的原子、桥梁、电磁场的一个模式——
都受到与位移成正比的回复力，因为每一个光滑势
在其势阱底部都是一条抛物线。因此，只需求解一次
量子谐振子，便求得整个世界微\omterm{prop:b3:lagrangian-mechanics:modes}{小振动}的解。本章以
量子力学标志性的代数方法求解它：两个\omterm{def:b3:harmonic-oscillator:ladder}{升降算符}
沿着一个完美均匀的能级阶梯上升和下降，
能级为 $\big(n +
\tfrac12\big)\hbar\omega$，底部的半步
——\omterm{thm:b3:harmonic-oscillator:spectrum}{零点能}——是一个定理，而非一种选择。其影响
范围广泛，从氦为何始终不凝固，到二氧化碳分子
以自身的 $\hbar\omega$ 振动时
如何截获地球向外发射的
热量。

\section{阶梯}

\begin{definition}[升降算符]\label{def:b3:harmonic-oscillator:ladder}
\index{升降算符}\index{湮灭算符}\index{产生算符}
对于 $\hat H = \hat p^2/2m + \tfrac12 m\omega^2\hat x^2$，引入
无量纲的非\omterm{def:b3:quantum-formalism:observable}{厄米算符}对
\[
\hat a = \sqrt{\frac{m\omega}{2\hbar}}\Big(\hat x +
\frac{\iu\hat p}{m\omega}\Big) , \qquad
\hat a^\dagger = \sqrt{\frac{m\omega}{2\hbar}}\Big(\hat x -
\frac{\iu\hat p}{m\omega}\Big) .
\]
由 $[\hat x, \hat p] = \iu\hbar$ 可得：
\[
[\hat a, \hat a^\dagger] = 1 , \qquad
\hat H = \hbar\omega\Big(\hat N + \tfrac12\Big) , \quad
\hat N = \hat a^\dagger\hat a .
\]
$\hat N$ 是\omterm{def:b3:quantum-formalism:observable}{厄米算符}；其\omterm{def:b3:quantum-formalism:observable}{本征值}将计数\emph{量子}，而
$\hat a$、$\hat a^\dagger$——\emph{湮灭}算符和
\emph{产生}算符——分别移除和增加一个量子。
\end{definition}

\begin{theorem}[仅由代数得到的谱]\label{thm:b3:harmonic-oscillator:spectrum}
\index{零点能}
$\hat H$ 的\omterm{def:b3:quantum-formalism:observable}{本征值}恰为
\[
E_n = \Big(n + \tfrac12\Big)\hbar\omega , \qquad n = 0, 1, 2, \dots
\]
——一串无穷的等距阶梯，步长为 $\hbar\omega$，位于一个
\emph{不为零}的基态能级之上，即\emph{\omterm{thm:b3:harmonic-oscillator:spectrum}{零点能}}
$\tfrac12\hbar\omega$。归一化本征态 $\ket n$ 通过
下列关系相联系：
\[
\hat a\ket n = \sqrt n\,\ket{n - 1} , \qquad
\hat a^\dagger\ket n = \sqrt{n + 1}\,\ket{n + 1} ,
\qquad
\ket n = \frac{(\hat a^\dagger)^n}{\sqrt{n!}}\,\ket0 .
\]
\end{theorem}

\begin{proof}
由\omterm{def:b3:quantum-formalism:commutator}{对易子}可得 $\hat N\hat a = \hat a(\hat N - 1)$：若
$\ket\nu$ 的 $\hat N$ \omterm{def:b3:quantum-formalism:observable}{本征值}为 $\nu$，则 $\hat a\ket\nu$ 是
\omterm{def:b3:quantum-formalism:observable}{本征值}为 $\nu - 1$ 的本征矢（或零矢量）。每一步下降仅在
$\nu \ge 0$ 时才被允许，因为 $\nu = \braket{\nu}{\hat
N\nu} = \|\hat a\ket\nu\|^2 \ge 0$；因此下降必须终止，而它只能终止于
满足 $\hat a\ket{\nu_0}
= 0$，由此 $\nu_0 = 0$。故 $\nu$ 取遍非负
整数。归一化系数由 $\|\hat a\ket n\|^2 = n$
及 $\|\hat a^\dagger\ket n\|^2 = n + 1$ 得出。
\end{proof}

\begin{omfigure}
\begin{tikzpicture}[font=\small, >=Stealth]
  \draw[->] (0,-0.3) -- (0,5.4);
  \node[anchor=south] at (0,5.4) {$E$};
  \foreach \n in {0,1,2,3,4}
    {\draw[very thick, omDef] (0.5,{0.55+\n*1.05}) -- (3.3,{0.55+\n*1.05});
     \node[anchor=west] at (3.45,{0.55+\n*1.05})
       {$\ket{\n}$\; $E = (\n + \tfrac12)\hbar\omega$};}
  \draw[->, omThm, very thick] (1.35,0.62) -- (1.35,1.53);
  \node[omThm, anchor=east] at (1.28,1.08) {$\hat a^\dagger$};
  \draw[->, omExo, very thick] (2.45,1.53) -- (2.45,0.62);
  \node[omExo, anchor=west] at (2.52,1.08) {$\hat a$};
  \draw[<->] (0.15,0) -- (0.15,0.55);
  \node[anchor=west, font=\footnotesize] at (0.28,0.22) {$\tfrac12\hbar\omega$};
  \draw[dashed] (0,0) -- (3.3,0);
\end{tikzpicture}

{\small 谐振子的阶梯：等距台阶 $\hbar\omega$，由
$\hat a^\dagger$ 向上攀登、由 $\hat a$ 向下回落，最低处立于
比经典\omterm{def:b3:relativistic-dynamics:momentum-energy}{静能}高半个台阶的地面上。}
\end{omfigure}

\begin{proposition}[空间中的态]\label{prop:b3:harmonic-oscillator:wavefunctions}
基态为高斯函数
\[
\varphi_0(x) = \Big(\frac{m\omega}{\pi\hbar}\Big)^{1/4}
\eu^{-m\omega x^2/2\hbar} ,
\]
它由求解\emph{一阶}方程 $\hat a\varphi_0
= 0$ 得到；它使海森堡不等式 $\Delta x\,\Delta p =
\hbar/2$ 取等号，其中 $\Delta x = \sqrt{\hbar/2m\omega}$。
反复作用 $\hat a^\dagger$ 即生成 $\varphi_n$：一个 $n$ 次多项式
（有 $n$ 个节点）乘以同一高斯函数。对于大的 $n$，
$|\varphi_n|^2$ 快速振荡于经典停留时间分布附近，
在转折点附近最大——
对应原理的一幅图像。
\end{proposition}

\begin{proof}[部分证明]
$\hat a\varphi_0 = 0$ 即 $\varphi_0' = -(m\omega/\hbar)x
\varphi_0$：归一化的高斯函数。饱和性：高斯函数是
\cref{thm:b3:quantum-formalism:uncertainty}中取等号的情形。
多项式结构源于 $\hat a^\dagger$ 对 $x$
和 $\dd/\dd x$ 均为一阶；大 $n$ 的结论在
\cref{exo:b3:harmonic-oscillator:12}中检验。
\end{proof}

\begin{omfigure}
\begin{tikzpicture}
  \begin{axis}[omaxis, axis x line=bottom, axis y line=left,
    width=6.6cm, height=5.2cm, xmin=-4, xmax=4, ymin=0, ymax=3.6,
    xlabel={$x$（单位 $\sqrt{\hbar/m\omega}$）}, ylabel={},
    xlabel style={at={(axis description cs:0.5,-0.1)}, anchor=north},
    xtick={-2,0,2}, ytick=\empty, clip=false]
    \addplot[omDef, thick, domain=-4:4, samples=90]
      {0.56*exp(-x*x)+0.25};
    \addplot[omThm, thick, domain=-4:4, samples=90]
      {1.13*x*x*exp(-x*x)+1.3};
    \addplot[omExo, thick, domain=-4:4, samples=120]
      {0.28*(2*x*x-1)^2*exp(-x*x)+2.35};
    \node[omDef, font=\scriptsize, anchor=west] at (axis cs:2.4,0.5) {$|\varphi_0|^2$};
    \node[omThm, font=\scriptsize, anchor=west] at (axis cs:2.4,1.55) {$|\varphi_1|^2$};
    \node[omExo, font=\scriptsize, anchor=west] at (axis cs:2.4,2.6) {$|\varphi_2|^2$};
  \end{axis}
  \begin{axis}[omaxis, axis x line=bottom, axis y line=left,
    width=6.6cm, height=5.2cm, xmin=-5.2, xmax=5.2, ymin=0, ymax=0.45,
    xshift=7.1cm,
    xlabel={$x$}, ylabel={},
    xlabel style={at={(axis description cs:0.5,-0.1)}, anchor=north},
    xtick=\empty, ytick=\empty, clip=false]
    \addplot[omDef, domain=-4.6:4.6, samples=300]
      {(64*x^6-480*x^4+720*x^2-120)^2*exp(-x*x)/81680};
    \addplot[omThm, very thick, dashed, domain=-3.54:3.54, samples=120]
      {0.3183/sqrt(max(0.35,13-x*x))};
    \node[omThm, font=\scriptsize, anchor=south] at (axis cs:0,0.36) {经典停留时间};
    \node[omDef, font=\scriptsize, anchor=west] at (axis cs:1.9,0.42) {$|\varphi_6|^2$};
  \end{axis}
\end{tikzpicture}

{\small 左：最低几个态的概率密度（竖直方向错开）：
无节点的高斯函数，其后依次为具有 $n$ 个节点的 $\varphi_n$。右：
在 $n$ 很大时，量子密度围绕经典
分布振荡；经典分布在转折点处堆积，
振动质点在此停留较久。}
\end{omfigure}

\section{零点能是真实的}

\begin{remark}[为何能量不能再低]\label{rem:b3:harmonic-oscillator:zeropoint}
能量为零的态需要 $\langle\hat p^2\rangle =
\langle\hat x^2\rangle = 0$：既完全静止，\emph{又}完全
居中，这被 $\Delta x\,\Delta p \ge \hbar/2$ 禁止。基
态是该不等式所允许的最佳折中，而
$\tfrac12\hbar\omega$ 即这笔“租金”。它绝非记账用的常数：
零点运动使 X 射线衍射图样在绝对零度时变模糊；
它使较轻的同位素的有效键较弱（H$_2$ 与
D$_2$ 在同一电子势阱中的解离能可测得不同）；
而在氦中，零点运动如此剧烈——原子轻、
吸引微弱，以至于液体\emph{从不}
在自身蒸气压下凝固：它是绝对零度时唯一
仍为液体的元素，仅在 25 个大气压的协助下才凝固。
\end{remark}

\begin{example}[$\hbar\omega$ 的量级]\label{ex:b3:harmonic-oscillator:scales}
一氧化碳键（$\omega = \qty{4.1e14}{rad/s}$）：$\hbar\omega =
\qty{0.27}{eV}$，为室温下 $k_{\text{B}}T$ 的十倍——
分子振动在日常空气中被冻结，这正是
双原子热容曾令十九世纪的人们困惑。单摆
（$\omega = \qty{5}{rad/s}$）：$\hbar\omega = \qty{3e-15}{eV}$，
完全超出分辨能力——正是对应极限。一面 LIGO
镜子，质量为 \qty{40}{kg}，悬挂成以 \qty{100}{Hz} 摆动，
作为一个振子模式，其零点振幅 $\Delta x =
\sqrt{\hbar/2m\omega} \approx \qty{5e-20}{m}$——而该观测站
例行地分辨\emph{恰在}这一量子下限处的位移：
四十千克物体的零点运动如今已成为工程
约束（见\cref{exo:b3:harmonic-oscillator:3}）。
\end{example}

\section{相干态：量子振子的经典一面}

\begin{definition}[相干态]\label{def:b3:harmonic-oscillator:coherent}
\index{相干态}
\emph{相干态} $\ket\alpha$ 是\omterm{def:b3:harmonic-oscillator:ladder}{湮灭算符}的
本征态，$\hat a\ket\alpha = \alpha\ket\alpha$，其中
$\alpha$ 为任意复数。在数态基下展开，
\[
\ket\alpha = \eu^{-|\alpha|^2/2}\sum_{n=0}^\infty
\frac{\alpha^n}{\sqrt{n!}}\,\ket n :
\]
量子计数服从泊松分布，均值为 $\langle\hat
N\rangle = |\alpha|^2$，涨落 $\Delta N = |\alpha|$。
\end{definition}

\begin{proposition}[为什么激光是经典的]\label{prop:b3:harmonic-oscillator:coherent-motion}
\omterm{def:b3:harmonic-oscillator:coherent}{相干态}即在\omterm{def:b3:hamiltonian-mechanics:hamiltonian}{相空间}中平移了的基态高斯函数；
在振子的演化下，它\emph{始终}保持相干，
$\alpha(t) = \alpha\,\eu^{-\iu\omega t}$：其中心精确地
做经典运动，$\langle\hat x\rangle(t) =
x_0\cos\omega t + \cdots$，而其宽度保持最小值
$\Delta x\,\Delta p = \hbar/2$ 永远成立——一个永不
扩散的波包。\omterm{def:b3:harmonic-oscillator:coherent}{相干态}是量子振子模拟经典振子的途径：
激光的光即某个场模式的\omterm{def:b3:harmonic-oscillator:coherent}{相干态}，
其光子数服从泊松分布（即每台光电探测器的
\emph{散粒噪声}），其场如麦克斯韦所述那样振荡。
\end{proposition}

\begin{proof}[部分证明]
将 $\ket\alpha = \sum c_n\ket n$ 代入
$\hat a\ket\alpha = \alpha\ket\alpha$，可得 $c_n = \alpha c_{n-1}/\sqrt
n$。演化时，每个 $\ket n$ 获得 $\eu^{-\iu(n + 1/2)\omega t}$，
重新求和后得到 $\alpha\eu^{-\iu\omega t}$ 的\omterm{def:b3:harmonic-oscillator:coherent}{相干态}，
（忽略整体相位）。$\langle\hat x\rangle \propto
\operatorname{Re}\alpha(t)$，由 $\hat x \propto \hat a +
\hat a^\dagger$ 可得。该态为平移后的高斯函数，
这一点在本阶段不加证明地承认。
\end{proof}

\begin{omfigure}
\begin{tikzpicture}[font=\small, >=Stealth]
  \draw[->] (-3.1,0) -- (3.4,0) node[right] {$\langle\hat x\rangle$};
  \draw[->] (0,-3.0) -- (0,3.2) node[above] {$\langle\hat p\rangle$};
  \draw[gray, dashed] (0,0) circle (2.2);
  \fill[omDef!25, draw=omDef, thick] (2.2,0) circle (0.42);
  \fill[omDef!25, draw=omDef, thick] (50:2.2) circle (0.42);
  \fill[omDef!25, draw=omDef, thick] (100:2.2) circle (0.42);
  \draw[->, thick] (2.53,0.9) arc (20:40:2.6);
  \node[anchor=west] at (2.62,1.35) {$\omega$};
  \fill[omThm!30, draw=omThm, thick] (0,0) circle (0.42);
  \node[omThm, font=\footnotesize, anchor=east] at (-0.3,-0.68) {$\ket0$};
  \node[omDef, font=\footnotesize, anchor=west] at (2.3,-0.62) {$\ket\alpha$};
  \node[font=\footnotesize, anchor=west] at (-3.05,2.75)
    {最小不确定度斑点 $\Delta x\,\Delta p = \hbar/2$，经典地绕行};
\end{tikzpicture}

{\small \omterm{def:b3:hamiltonian-mechanics:hamiltonian}{相空间}图像（比较
\cref{ch:b3:hamiltonian-mechanics}）：基态是原点处的最小
不确定度斑点；\omterm{def:b3:harmonic-oscillator:coherent}{相干态}即同一斑点
经平移后以 $\omega$ 绕行而不变形——量子
力学对经典振动的最佳模仿。}
\end{omfigure}

\begin{method}[振子代数]\label{met:b3:harmonic-oscillator:method}
(1) 将所求量用 $\hat a$、$\hat a^\dagger$ 表示：$\hat x
= \sqrt{\hbar/2m\omega}\,(\hat a + \hat a^\dagger)$, $\hat p =
\iu\sqrt{m\hbar\omega/2}\,(\hat a^\dagger - \hat a)$。(2) 把 $\hat a$ 移到右侧，
利用 $[\hat a, \hat a^\dagger] = 1$；使用
$\hat a\ket0 = 0$。(3) 矩阵元：$\hat x$ 只连接
相邻能级——振动光谱学中的选择定则 $\Delta n = \pm1$。
(4) 任何二次\omterm{def:b3:hamiltonian-mechanics:hamiltonian}{哈密顿量}（耦合振子、电路、场模式）
均可对角化为互相独立的
阶梯：先求\omterm{prop:b3:lagrangian-mechanics:modes}{简正模}，再逐一量子化。(5) 先看数值：
$\hbar\omega$ 与 $k_{\text{B}}T$ 的比较决定
体系是量子的还是经典的，应先于一切代数运算。
\end{method}

\section{练习}

\begin{exercise}[$\star$]\label{exo:b3:harmonic-oscillator:1}
(a) 验证 $[\hat a, \hat a^\dagger] = 1$，由 $[\hat x, \hat p] =
\iu\hbar$ 出发。(b) 直接展开验证 $\hat H = \hbar\omega(\hat a^\dagger\hat a +
\tfrac12)$。(c) 求逆以表示 $\hat x$ 与
$\hat p$。(d) 证明 $\hat N$ 是厄米的，并说明为何 $\hat a$
单独不能作为可观测量。
\end{exercise}

\begin{exercise}[$\star$]\label{exo:b3:harmonic-oscillator:2}
在本征态 $\ket n$ 上：(a) 证明 $\langle\hat x\rangle =
\langle\hat p\rangle = 0$；(b) 计算 $\langle\hat x^2\rangle$ 与
$\langle\hat p^2\rangle$；(c) 推出 $\Delta x\,\Delta p = (n +
\tfrac12)\hbar$；(d) 检验维里比 $\langle E_k\rangle =
\langle E_p\rangle$。
\end{exercise}

\begin{exercise}[$\star$]\label{exo:b3:harmonic-oscillator:3}
零点振幅 $\Delta x = \sqrt{\hbar/2m\omega}$：计算
分别求 (a) CO 分子（$\mu = \qty{1.14e-26}{kg}$，$\omega =
\qty{4.1e14}{rad/s}$），与键长
\qty{0.11}{nm} 相比较；(b) 固体中的氢原子（$m =
\qty{1.7e-27}{kg}$，$\hbar\omega = \qty{0.1}{eV}$）；(c) LIGO 的镜子
（$m = \qty{40}{kg}$，$f = \qty{100}{Hz}$）；(d) 将三者
按各自体系尺度的分数排序，并评论谁算是
“量子”。
\end{exercise}

\begin{exercise}[$\star$]\label{exo:b3:harmonic-oscillator:4}
(a) 利用升降关系，计算 $\bra m\hat x\ket n$，并
证明除非 $m = n \pm 1$，否则其为零。(b) 推出
光吸收的振动选择定则 $\Delta n = \pm1$（耦合项
为 $\propto\hat x$）。(c) 为何异核分子（CO）
吸收红外线，而 N$_2$ 却不吸收？(d) 实际分子会出现
位于 $\approx 2\hbar\omega$ 附近的较弱“泛频”谱线，
这说明势函数的什么性质？
\end{exercise}

\begin{exercise}[$\star\star$]\label{exo:b3:harmonic-oscillator:5}
(a) 将 $\hat a\varphi_0 = 0$ 视为微分方程求解，并
归一化。(b) 生成 $\varphi_1$ 与 $\varphi_2$，
方法是作用 $\hat a^\dagger$。(c) 仅由宇称即可验证 $\varphi_1 \perp \varphi_0$。
(d) 画出三者的概率密度示意图，并检验节点数。
\end{exercise}

\begin{exercise}[$\star\star$]\label{exo:b3:harmonic-oscillator:6}
玻尔兹曼预览。一组相同的振子
在温度 $T$ 下，振子占据能级 $n$ 的概率 $\propto
\eu^{-E_n/k_{\text{B}}T}$（参见本科二年级分册中的玻尔兹曼因子）。(a)
证明平均量子数为 $\langle n\rangle =
1/(\eu^{\hbar\omega/k_{\text{B}}T} - 1)$。(b) 计算 CO 振动在
\qty{300}{K} 与 \qty{2000}{K} 下的值。(c) 证明平均
能量在高温时趋于 $k_{\text{B}}T$（恢复能量均分
定理），低温时趋于 $\tfrac12\hbar\omega$。(d) 在何温度下
\qty{15}{\micro m} 振动（二氧化碳的
弯曲振动）平均量子数恰为 $\langle n\rangle = 0.1$？
\end{exercise}

\begin{exercise}[$\star\star$]\label{exo:b3:harmonic-oscillator:7}
\omterm{def:b3:harmonic-oscillator:coherent}{相干态}统计。(a) 由 $\ket\alpha$ 的展开式，
证明 $\mathcal P(n)$ 为均值 $|\alpha|^2$ 的泊松分布。(b) 证明
$\langle\hat N\rangle = |\alpha|^2$ 与 $\Delta N = |\alpha|$。(c) 一台
\qty{1}{mW} 的激光器，波长为 \qty{633}{nm}：每秒光子数，以及
一秒内的相对涨落 $\Delta N/\langle N\rangle$。(d)
散粒噪声：证明光电流的噪信比以
$1/\sqrt{\langle N\rangle}$ 的形式下降——为何明亮光束看起来平滑。
\end{exercise}

\begin{exercise}[$\star\star$]\label{exo:b3:harmonic-oscillator:8}
\omterm{def:b3:harmonic-oscillator:coherent}{相干态}的演化。(a) 将演化相位乘到
展开式上，并证明 $\ket{\alpha(t)}$，其中 $\alpha(t) =
\alpha\eu^{-\iu\omega t}$（相差一个整体相位）。(b) 推出
$\langle\hat x\rangle(t)$ 与 $\langle\hat p\rangle(t)$，并与
经典解比较。(c) 为何能量本征态尽管是定态，
\emph{不}描述摆动的单摆
——\omterm{def:b3:harmonic-oscillator:coherent}{相干态}的哪一特征能做到？(d) 估计 $|\alpha|$
对实际单摆（\qty{10}{g}，\qty{10}{cm} 振幅，
\qty{1}{Hz}），并加以评论。
\end{exercise}

\begin{exercise}[$\star\star$]\label{exo:b3:harmonic-oscillator:9}
量子 LC 电路。一个超导回路，其 $L =
\qty{10}{nH}$、$C = \qty{0.4}{pF}$，电荷与磁通量
如 $x$ 与 $p$ 那样振荡。(a) 其共振频率（本科一年级分册）及
量子能量 $\hbar\omega$（以 $\unit{\micro eV}$ 计）。(b) 在何温度以下
$k_{\text{B}}T \ll \hbar\omega$ 时，电路
处于基态？(c) 超导量子比特为何要在
\qty{20}{mK} 的稀释制冷机中运行？(d) 零点电压
涨落 $\Delta V = \Delta q/C$，其中 $\Delta q =
\sqrt{\hbar\omega C/2}$，求其数值。
\end{exercise}

\begin{exercise}[$\star\star\star$]\label{exo:b3:harmonic-oscillator:10}
零点运动导致的范德瓦耳斯力。将两个中性原子
（相距 $R$）视为两个相同的偶极振子（电荷 $e$，质量
$m$、频率 $\omega_0$），其位移通过偶极能
$\lambda\,x_1x_2$ 耦合，其中 $\lambda = e^2/2\pi\varepsilon_0R^3$。
(a) 证明简正坐标 $(x_1 \pm x_2)/\sqrt2$ 以
$\omega_\pm = \omega_0\sqrt{1 \pm \lambda/m\omega_0^2}$。(b) 基态
能为 $\tfrac\hbar2(\omega_+ + \omega_-)$，展开到
$\lambda$ 的二阶，并证明相互作用能为
\[
\Delta E = -\frac{\hbar\lambda^2}{8m^2\omega_0^3}
\ \propto\ -\frac{1}{R^6} .
\]
(c) 为何这种吸引具有普遍性——原子间即使
完全没有永久偶极矩也存在？(d) $1/R^6$ 律就是本科一年级分册中
真实气体修正所对应的范德瓦耳斯力：从微观看，
“涨落”究竟指什么——以及
这一力在 $\hbar = 0$ 的世界中会如何？
\end{exercise}

\begin{exercise}[$\star\star\star$]\label{exo:b3:harmonic-oscillator:11}
囚禁离子的边带。位于谐振势阱中的离子（$f =
\qty{1}{MHz}$）吸收激光。由于离子运动，其
吸收光谱在 $\nu_0$ 处显示电子谱线，两侧伴有
位于 $\nu_0 \pm f$ 处的谱线：这些跃迁使\emph{运动}
量子数改变 $\mp1$，并伴随电子态的改变。(a) 求
$\hbar\omega_{\text{阱}}$ 以 $\unit{neV}$ 为单位的值，并说明分辨
边带为何需要极窄的谱线？(b) 驱动 $\nu_0 - f$ 谱线
每个循环移除一个运动量子，请解释\emph{边带
冷却}使离子达到基态。(c) 一旦 $\langle n\rangle \approx
0$，下边带完全消失——这种
不对称性为何能证明已达到量子基态？(d) 这正是
确认单个原子运动基态的方式——而对千克量级的
LIGO 镜子则借助其他手段——请说明
离子所达到的“温度”，取 $f = \qty{1}{MHz}$ 与
$\langle n\rangle = 0.05$ 时的值。
\end{exercise}

\begin{exercise}[$\star\star\star$]\label{exo:b3:harmonic-oscillator:12}
定量地看经典极限。一个振幅为
$A$ 的经典振子在 $[x, x + \dd x]$ 内所占的比例为 $\dd
t/T = \dd x/\pi\sqrt{A^2 - x^2}$。(a) 推导该停留时间
分布。(b) 对量子态 $n$，取 $A_n$，使 $E_n =
\tfrac12 m\omega^2A_n^2$，并将经典分布与
（已给出的）局域平均 $|\varphi_n|^2$ 相比较：两者何处一致，
何处必然不同？(c) 证明相邻能级之间的
相对间距 $\Delta E/E$ 随 $1/n$ 趋于零：能量变得
实际上连续。(d) 对
\cref{exo:b3:harmonic-oscillator:8}(d) 中的单摆，估计 $n$ 与 $\Delta
E/E$，并用一句话概括对应论证的结论。
\end{exercise}

\section{问题：使地球变暖的分子}

\begin{problem}[{周末问题——二氧化碳的量子阶梯
与温室效应}]\label{pb:b3:harmonic-oscillator:1}
二氧化碳分子是线性的 O=C=O 链，由若干个量子化的
振子组成。其弯曲模的 $\hbar\omega$ 恰好落在
地球向外发出的热辐射中部，这正是这种痕量气体
——每万个分子中约四个——引导着行星的
气候。数据：弯曲模波数 $\tilde\nu_2 = \qty{667}{cm^{-1}}$
（光谱学家所用的单位：$E = hc\tilde\nu$，$\qty{1}{cm^{-1}} =
\qty{1.24e-4}{eV}$）；不对称伸缩 $\tilde\nu_3 =
\qty{2349}{cm^{-1}}$；对称伸缩 $\tilde\nu_1 =
\qty{1388}{cm^{-1}}$；$k_{\text{B}}T$ 在 $\qty{288}{K}$ 时为
\qty{24.8}{meV}；维恩位移定律（本科二年级分册）：$\lambda_{\max}T =
\qty{2898}{\micro m.K}$。

\medskip\noindent\textbf{第一部分——线性分子的振动模式。}
\begin{enumerate}
  \item 三个原子具有九个\omterm{def:b3:lagrangian-mechanics:coordinates}{自由度}：其中几个为整个分子的平动，
        几个为转动（分子是
        \emph{线性}的），剩下几个为振动？
  \item 描述四种振动：对称伸缩、不对称
        伸缩，以及二重简并的弯曲——为何弯曲模
        出现两次？
  \item 光与振动的\emph{电偶极子}相耦合
        （见\cref{exo:b3:harmonic-oscillator:4}）。O=C=O 的哪些振动模式
        调制偶极矩，哪一个是
        红外不活跃的？
  \item 将三个波数换算为光子波长，并
        标出它们的位置：哪些位于热红外区？
  \item 为何空气中两种主要气体 N$_2$ 与 O$_2$ 几乎完全
        不吸收红外线——这对大气的
        透明度有何影响？
  \item 水蒸气是弯曲的偶极分子，在红外的大部分
        波段都有吸收：在哪一光谱“窗口”中，二氧化碳的弯曲模
        主要单独起作用（比较 \qty{15}{\micro m} 与
        水的强吸收带位于 \qty{8}{\micro m} 以下及
        \qty{20}{\micro m} 以上）？
\end{enumerate}

\medskip\noindent\textbf{第二部分——弯曲模的量子阶梯。}
\begin{enumerate}[resume]
  \item 以 meV 为单位计算 $\hbar\omega_2$，并求出能级阶梯 $E_n$。
  \item 占据 $n = 1$ 的分子比例为多少，
        在 \qty{288}{K} 下（相对 $n = 0$，按玻尔兹曼因子计）？
        $n = 2$ 呢？
  \item 哪一光子波长驱动 $n = 0 \to 1$？为何
        \emph{同一}波长会主导发射？
  \item 由谐振子阶梯证明，同一波长适用于
        整个阶梯（$n \to n + 1$，对每个 $n$）：起作用的
        是能级间距的何种性质？
  \item 分子还会转动，并带来精细结构：
        \qty{15}{\micro m} 处的特征实际上是一个\emph{谱带}，宽约
        \qty{1}{\micro m}。对温室气体而言，谱带\emph{宽度}
        很重要（想一想，一旦谱带中心
        变得不透明会发生什么）？
  \item 空气中振动激发的 CO$_2$ 远更可能
        通过碰撞而非辐射失去其量子（辐射
        寿命 $\sim\qty{1}{s}$，碰撞时间 $\sim\qty{e-9}{s}$）：
        吸收的辐射能究竟去往何处？
  \item 反之，\qty{288}{K} 的空气中保有
        处于 $n = 1$ 的热布居（第 8 问）：该布居做了什么
        对能量收支有何重要影响？
\end{enumerate}

\medskip\noindent\textbf{第三部分——行星的辐射收支。}
\begin{enumerate}[resume]
  \item 将太阳视为温度 \qty{5800}{K} 的黑体，计算其
        维恩峰值。CO$_2$ 的弯曲模式能否截获较多太阳光？
  \item 将地表视为温度 \qty{288}{K} 的黑体，计算其
        维恩峰值，并在该热辐射曲线上标出 \qty{15}{\micro m}
        所对应的位置。
  \item 用四句话解释温室机制：
        太阳光射入，热红外射出，在 \qty{15}{\micro m} 处被截获，
        再向上\emph{与}向下两个方向辐射。
  \item 最终逸散入太空的再辐射来自高处的寒冷层
        （$\sim\qty{220}{K}$）：为何从更冷的层发射会
        使行星在这些波长上的出射功率减小
        （回想热辐射随 $T$ 增大而增强）？
  \item 增加 CO$_2$ 会使发射层升得更高、温度更低：
        请用一句话说明，此后地表为何必须升温，
        以重新达成收支平衡。
  \item 带中心已经不透明；新增
        CO$_2$ 的作用体现在谱带的\emph{翼部}，
        使辐射强迫随浓度呈对数增长：请将此
        与第 11 题联系起来。
\end{enumerate}

\medskip\noindent\textbf{第四部分——同位素：量子
指纹。}
\begin{enumerate}[resume]
  \item 振子频率与 $\sqrt{k/\mu}$ 成正比：对于不对称伸缩模式，
        将 $^{12}$C 替换为 $^{13}$C 会改变有效质量；
        观测到的谱线由
        \qty{2349}{cm^{-1}} 移至约 \qty{2283}{cm^{-1}}。请
        根据质量核对这一 $\approx 3\%$ 偏移的
        数量级。
  \item 调谐到这两条谱线的激光器可在气样中分别统计 $^{13}$CO$_2$ 与
        $^{12}$CO$_2$：请解释量子化的能级阶梯
        为何使得这种按同位素分辨的探测
        才有可能实现。
  \item 植物偏好较轻的同位素，故化石碳
        $^{13}$C 偏少。大气 CO$_2$ 的实测同位素比值随其
        浓度升高发生了何种变化，并
        由此对新增气体的\emph{来源}可以
        得出什么结论？
  \item 同样的 \qty{15}{\micro m} 物理过程也作用于金星
        （$96\%$ CO$_2$，\qty{90}{bar}）：其
        \qty{737}{K} 地表说明了什么？
  \item 火星大气同样几乎全为 CO$_2$，但气压仅为 \qty{6}{mbar}，
        且极为寒冷：金星–地球–火星这一组对比
        能说明哪一变量控制着该效应的
        强弱？
  \item 概括这一著名结果：一个 \qty{83}{meV} 的量子——
        即弯曲振动三原子分子的 $\hbar\omega$——，停留在
        一颗 \qty{288}{K} 行星热谱上的 \qty{15}{\micro m} 处，
        被吸收、在纳秒内热化，并
        从寒冷高空再辐射出去，这一机制使得
        空气中占 $0.04\%$ 的成分决定了地球的温度。
\end{enumerate}
\end{problem}
